\documentclass[12pt, a4paper]{article}
\usepackage[utf8]{inputenc}
\usepackage[spanish]{babel}
\usepackage{graphicx}
\usepackage{hyperref}
\usepackage{amsmath}
\usepackage{listings}
\usepackage{xcolor}
\usepackage{geometry}
\usepackage{enumitem}
\usepackage{titlesec}
\usepackage{booktabs}
\usepackage{longtable}
\usepackage{array}
\usepackage{tabularx}
\usepackage{tikz}
\usepackage{pdfpages}
\usepackage{float}
\usepackage{caption}
\usepackage{subcaption}
\usepackage{microtype}
\graphicspath{{media/}}
\usetikzlibrary{positioning}
\geometry{margin=2.5cm}
\setlength{\parskip}{0.6em}
\setlength{\parindent}{0pt}

\newcommand{\nodo}[1]{\fbox{\parbox{3cm}{\centering #1}}}
\newcommand{\flecha}{\par\bigskip\centering $\downarrow$\par\bigskip}

% Configuración de secciones
\titleformat{\section}{\large\bfseries}{\thesection}{1em}{}
\titleformat{\subsection}{\normalsize\bfseries}{\thesubsection}{1em}{}


\begin{document}

% ===== I. PRESENTACIÓN =====
\begin{titlepage}
\centering
\textbf{Universidad de la Habana \\
Facultad de Matemática y Computación \\
Carrera: Ciencia de la Computación \\
Año: 3ero \\
Asignaturas: Ingeniería de Software y Sistemas de Bases de Datos II}\\[2cm]


\begin{figure}[H]
	\centering
	
	% --- Subfigura (a) ---
	\begin{subfigure}[b]{0.45\textwidth}
		\centering
		\includegraphics[width=0.9\textwidth]{logo.png}
		\caption{Logo}
		\label{fig:logo}
	\end{subfigure}
	\hfill
	% --- Subfigura (b) ---
	\begin{subfigure}[b]{0.45\textwidth}
		\centering
		\includegraphics[width=0.9\textwidth]{logoMin.jpeg}
		\caption{Miniatura}
		\label{fig:logoMin}
	\end{subfigure}
	
	% --- Leyenda general de la figura ---
	\caption{Arte de la aplicación}
	\label{fig:art}

\end{figure}


\vspace{1cm}
{\LARGE \textbf{Sistema de Gestión de Misiones de Jujutsu Kaisen}}\\[0.5cm]
{\large \textbf{Informe Técnico}}

\vspace{1.5cm}
\textbf{Desarrolladores:}\\
\begin{tabular}{lll}
\textbf{Nombre y Apellidos} & \textbf{Grupo} & \textbf{Usuario Telegram} \\
Ronald Alfonso Pérez & C-311 & (\href{https://t.me/RAlfonsoP}{\textcolor{blue}{@RAlfonsoP}}),\\
Juan Carlos Carmenate Díaz & C-311 & (\href{https://t.me/Juank404}{\textcolor{blue}{@Juank404}}),\\
Sebastian González Alfonso & C-312 & (\href{https://t.me/sebagonz106}{\textcolor{blue}{@sebagonz106}}),\\
Abraham Rey Sánchez Amador & C-312 & (\href{https://t.me/AbrahamR_Sanchez}{\textcolor{blue}{@AbrahamR\_Sanchez}}).
\end{tabular}

\vfill
Curso: 2025-2026
\end{titlepage}

\begin{abstract}
	
\end{abstract}

\tableofcontents
\newpage

% ===== REFERENCIAS =====
\begin{thebibliography}{9}
\bibitem{lucina}
García H, L. \& Montes de Oca R, M. (2005). \textit{Sistemas de Bases de Datos: Modelación y Diseño - Capítulos para estudiantes}. Versión digital.

\bibitem{pressman2010}
Pressman S. R. (2010). \textit{Ingeniería del Software. Un enfoque práctico} (Séptima edición). McGraw-Hill Interamericana Editores, México, D. F.

\bibitem{fowler2002}
Fowler, M. (2002). \textit{Patterns of Enterprise Application Architecture}. Addison-Wesley.
\end{thebibliography}

\end{document}